\chapter*{Abstract}

The study of this undergraduate thesis aims to analyze the code produced by Large Language Models in solving senior-level problems from the Brazilian Olympiad in Informatics, Programming Modality. Due to this, experiments were conducted with five different languages, the use of images and descriptions, providing time and memory constraints, and, finally, prompt techniques.
For this purpose, a tool developed in Python was created with the objective of automating code generation through model API requests, verifying whether the responses produced by the models contain executable code for the Judge to compile and execute in the machine's own environment.
Furthermore, all requests to the LLMs were made with a temperature equal to zero, maintaining the characteristic of greedy decoding, that is, reducing the variation of the codes produced by the models. This directly impacts the choice of the $\passat{1}$ metric ($\passat{k}$, such that $k = 1$) for the Judge's verdicts, since the objective of the study is not to evaluate the ``creativity'' of the models in solving senior-level problems, but rather the ability to solve the same problems in different scenarios.
Thus, the study shows that, in this dataset, the syntactic complexity of the languages has a direct impact on problem resolution; moreover, even when using the best language (C++), which performed best on average across the models, the use of Base64-formatted images in the prompts had a direct impact on token consumption and does not guarantee problem resolution. However, in the language that had the worst performance (C), even when providing the execution time limit and memory limit, the produced solutions did not improve the responses reproduced by the LLMs. In addition, the use of prompt techniques helped the models improve their performance on $\passat{1}$ by providing a few simple input and output examples. Finally, it shows that large language models still face major challenges in producing solutions to programming problems, because the logic required to solve the question is the same, but the prediction of the next token has a direct impact, especially in languages that require implementations of data structures, optimizations, etc.

\vspace{5mm}
\noindent\textbf{\textit{Keywords}:~}Large Language Models; Competitive Programming; Brazilian Olympiad in Informatics; Experimental Evaluation; Prompt Engineering; Multimodal Vision.
